% ECONOMICS_PROBLEM: {"id": "L41-333", "title": "Dynamic competition, innovation and mergers", "jel_code": "L41", "source_id": "OP-0072"}
% !TeX program = xelatex
\documentclass[11pt,a4paper]{article}
\usepackage[margin=23mm]{geometry}
\usepackage{fontspec}
\setmainfont{Latin Modern Roman}
\usepackage{amsmath,amssymb,mathtools}
\usepackage{xcolor,enumitem,booktabs,longtable,array,xurl,hyperref}
\definecolor{muted}{HTML}{53636C}
\hypersetup{colorlinks=true,urlcolor=blue,linkcolor=blue}
\setlength{\parindent}{0pt}
\setlength{\parskip}{5pt}
\newcommand{\R}{\mathbb R}
\newcommand{\E}{\mathbb E}
\newcommand{\N}{\mathbb N}
\newcommand{\ind}{\mathbf 1}
\newcommand{\Var}{\operatorname{Var}}
\newcommand{\Cov}{\operatorname{Cov}}
\newcommand{\supp}{\operatorname{supp}}
\DeclareMathOperator*{\argmin}{arg\,min}
\DeclareMathOperator*{\argmax}{arg\,max}
\newcommand{\entrymeta}[3]{{\sffamily\small\color{muted}\textbf{\texttt{#1}}\quad|\quad #2\par #3\par}}
\newcommand{\Problemref}[1]{\texttt{#1}}
\newcommand{\JELref}[2]{\texttt{#2} (JEL \texttt{#1})}
\begin{document}
\section*{Dynamic competition, innovation and mergers}
\textbf{Problem ID:} \texttt{L41-333}\quad \textbf{JEL:} \texttt{L41}\par
% BEGIN ECONOMICS STATEMENT
\label{problem:OP-0072}
\label{atlas:prob:I2}

\noindent\textbf{Original identifier:} I2 (atlas item 72). \textbf{Classification:} Dynamic oligopoly and innovation research program. \textbf{Original tractability label:} Hard (editorial, not a complexity result).

\subsection*{Definitions and assumptions}

Let firms choose research investment $a_i$, entry or exit, acquisition offers, and product actions in a discounted stochastic game. The public state $x$ records technologies, ownership, demand, and relevant private-information beliefs; discount factors satisfy $0<\delta_i<1$. A transition kernel $K(x'\mid x,a,\mu)$ depends on actions and merger policy $\mu$. If private information is retained, include each firm's posterior beliefs in its information state and specify Bayes-consistent updating and a corresponding Markov-perfect Bayesian equilibrium. Firm payoffs include acquisition payments and subtract investment and transaction resource costs; social welfare excludes pure payment transfers. A Markov-perfect equilibrium consists of strategies optimal after every state given rivals' strategies and the kernel. Let $\mathcal E(\mu)$ be the equilibrium set and fix or bound a policy-dependent selection rule $\sigma$. Define innovation using specified quality-adjusted products or productivity, not patent counts alone.

\subsection*{Formal research question}

For a baseline state distribution $F_0$, identify or bound the policy effects
\[
\Delta_I^{\sigma}(\mu,\mu_0)=\mathbb E_{F_0}\!\left[\sum_{t=0}^{T}\rho^t\{I_t(\mu,\sigma)-I_t(\mu_0,\sigma)\}\right],
\]
and the analogous discounted welfare difference, where $\rho\in(0,1]$ and $I_t$ is the declared innovation measure. Characterize sufficient conditions for the sign of the innovation effect to be invariant across admissible equilibria and parameter values. Jointly model product entry, acquisitions, and research direction, and assess predictions using exogenous policy or technology variation where available.

\subsection*{Interpretation and scope}

Concentration is an endogenous market outcome, so the causal effect of a merger policy cannot generally be identified by regressing innovation on concentration. An acquisition can internalize competition, change financing, acquire complementary technology, or remove a potential entrant; these mechanisms need distinct state transitions and costs. The welfare criterion should include consumer gains from new products, research resource costs, and displaced products, with a stated treatment of knowledge spillovers. Patent counts provide an imperfect observation equation and can change because of filing incentives. Counterfactual technologies that never appear require modeling assumptions and cannot be declared observed. Dynamic identification may fail when alternative investment costs and continuation values produce the same behavior; report sensitivity sets rather than one assumed model. Equilibrium selection can reverse comparative statics even when every estimated payoff is held fixed. A valuable extension would locate verifiable restrictions yielding sign-robust conclusions in a specified industry and horizon, instead of asserting one relationship for every market.

\subsection*{Source audit and status}

[S1] is historical motivation; its 1942 original U.S. edition is distinguished from the linked later publisher reprint. [S2] supplies a dynamic-industry framework and [S3] quantifies creative destruction in hard drives. None implies that concentration universally stimulates or suppresses research. These sources answer particular model and industry questions. The acquisition--innovation counterfactual and selection-robust sign target above are editorial research extensions, without a certified universal open status in 2026.

\subsection*{References}

\begin{enumerate}[label={[S\arabic*]},leftmargin=*]

\item Joseph A. Schumpeter (1942). \emph{Capitalism, Socialism and Democracy}. Original U.S. edition, Harper \& Brothers; publisher link is a later Routledge reprint. \url{https://www.routledge.com/Capitalism-Socialism-and-Democracy/Schumpeter/p/book/9780415107624}\par \textit{Source role:} Historical motivation, not a formal merger theorem. \textit{Locator:} Chapters 7--8; reprint contents identify first-edition preface as 1942.

\item Richard Ericson and Ariel Pakes (1995). \emph{Markov-Perfect Industry Dynamics: A Framework for Empirical Work}. Review of Economic Studies 62(1), 53--82. \url{https://doi.org/10.2307/2297841}\par \textit{Source role:} Dynamic oligopoly framework. \textit{Locator:} Abstract and article metadata.

\item Mitsuru Igami (2017). \emph{Estimating the Innovator's Dilemma: Structural Analysis of Creative Destruction in the Hard Disk Drive Industry, 1981--1998}. Journal of Political Economy 125(3), 798--847. \url{https://doi.org/10.1086/691524}\par \textit{Source role:} Industry-specific structural evidence. \textit{Locator:} Abstract and article metadata.

\end{enumerate}

\subsection*{Common conventions: Source mathematical conventions}

Unless an entry states otherwise, agent and firm sets are finite. State and action spaces are measurable, strategies and outcome maps are measurable, probability laws are specified, and expectations exist for the targets under discussion. Infinite-horizon discount factors lie in $(0,1)$ and discounted objectives are finite. Norms, tolerances and complexity input sizes must be specified for computational comparisons. Constants or primitives that appear locally are inputs, not empirical facts asserted by this catalogue.

\textbf{Identification.} A statistical model $m\in\mathcal M$ implies an observable law $P_m$ and target $\theta(m)$. At law $P$, its identified set is
\[
\Theta(P;\mathcal M)=\{\theta(m):m\in\mathcal M,\ P_m=P\}.
\]
Point identification holds when this set is a singleton, or equivalently when $P_{m_1}=P_{m_2}$ implies $\theta(m_1)=\theta(m_2)$ for all compatible models. Sharp bounds mean bounds attained, or approached when the boundary is not attained, by compatible models. Writing this set is a definition, not a solution: the substantive task is to establish informative restrictions and derive or estimate the set. An unrestricted-confounding model may give only trivial bounds.

\textbf{Inference.} Identification, estimability and valid uncertainty quantification are separate questions. Fix $\alpha\in(0,1)$. A confidence procedure $C_n$ over a declared law class $\mathcal P$ has asymptotic uniform coverage at level $1-\alpha$ when
\[
\liminf_{n\to\infty}\inf_{P\in\mathcal P}\Pr_P\{\theta(P)\in C_n\}\geq1-\alpha.
\]
For set-identified targets the coverage event must be defined separately, for example coverage of the full identified set. If an entry uses identification-indexed models instead of uniquely indexed laws, the same coverage convention is applied over those models. Pointwise limits do not establish uniform validity, and asymptotic validity does not guarantee finite-sample precision. Sample sizes, dependence structures and weak-identification sequences are part of each application.

\textbf{Counterfactuals.} $Y_i(a)$ is a potential outcome under a specified intervention, including its timing, implementation and versions. It is not $Y_i$ conditional on voluntarily choosing $a$. A full intervention vector permits interference; shorthand scalar treatment notation does not silently impose no interference. Assignment, selection, attrition, anticipation and measurement must be specified. A claim about an instrument's local effect is not automatically a population policy effect.

\textbf{Economic equilibrium.} A counterfactual model includes best responses, constraints, feasibility, market clearing, state transitions and expectations consistent with the stated information. When several equilibria exist, either a declared selector is an input or the target is set-valued. Comparisons across policy regimes must use the stated selection convention. Reduced-form relationships are not presumed invariant after an equilibrium-changing intervention.

\textbf{Optimization and welfare.} Optimization is over the stated feasible set. If attainment is not established, $\sup$ or $\inf$ and $\varepsilon$-optimal choices replace an asserted maximizer or minimizer. For example, with value $V=\sup_{a\in\mathcal A}W(a)<\infty$, an $\varepsilon$-optimal set is $\{a\in\mathcal A:W(a)\geq V-\varepsilon\}$ for $\varepsilon>0$. Benefits, costs, risk and health or environmental outcomes can be aggregated only using declared units and conversion weights. Interpersonal utility weights, the treatment of fiscal transfers and foreign profits, discounting, and the behavioral welfare criterion are explicit inputs. A comparative-static derivative additionally requires existence and appropriate differentiability of the selected counterfactual.

\textbf{Sources.} Local labels [S1], [S2], and so forth restart in every question. Locators identify the relevant abstract, model, section or result at the level actually checked; a bibliographic or abstract-level verification is not represented as inspection of an entire proof. Working papers are distinguished from later publications and changing versions. Source summaries are paraphrased. All URL checks and editorial formulations refer to the audit date stated above.
% END ECONOMICS STATEMENT
\end{document}
